\chapter{Gereksinimler}

\section{İşlevsel Gereksinimler (Functional Requirements)}

\subsection{Kullanıcı Yönetimi Gereksinimleri}
\textbf{FR-1.1: Kullanıcı Kayıt}
\begin{itemize}
    \item Sistem, yeni kullanıcıların kayıt olmasına izin vermelidir.
    \item Kayıt sırasında e-posta, kullanıcı adı ve şifre bilgileri toplanmalıdır.
    \item E-posta adresinin benzersizliği kontrol edilmelidir.
    \item Şifre en az 6 karakter olmalıdır.
    \item Kayıt sonrası JWT token döndürülmelidir.
\end{itemize}

\textbf{FR-1.2: Kullanıcı Doğrulama ve Giriş}
\begin{itemize}
    \item Sistem, kullanıcıların e-posta ve şifre ile giriş yapmasına izin vermelidir.
    \item Giriş başarılı olduğunda JWT token oluşturulmalıdır.
    \item Token'ın geçerlilik süresi belirlenmelidir.
    \item Giriş başarısız olduğunda uygun hata mesajı döndürülmelidir.
\end{itemize}

\textbf{FR-1.3: Profil Yönetimi}
\begin{itemize}
    \item Kullanıcılar profil bilgilerini (ad, soyad, profil resmi, durum mesajı) görüntüleyebilmelidir.
    \item Kullanıcılar profil bilgilerini güncelleyebilmelidir.
\end{itemize}

\subsection{Sohbet Yönetimi Gereksinimleri}
\textbf{FR-2.1: Birebir Sohbet}
\begin{itemize}
    \item Kullanıcılar, diğer kullanıcılarla birebir sohbet başlatabilmelidir.
    \item Sistem, aynı iki kullanıcı arasında mevcut sohbet varsa yeni sohbet oluşturmamalıdır.
\end{itemize}

\textbf{FR-2.2: Grup Sohbeti}
\begin{itemize}
    \item Kullanıcılar, birden fazla kullanıcı ile grup sohbeti oluşturabilmelidir.
    \item Grup sohbetine grup adı verilebilmelidir.
    \item Grup sohbetine katılımcı eklenebilmelidir.
    \item Grup sohbetinden katılımcı çıkarılabilmelidir.
\end{itemize}

\textbf{FR-2.3: Sohbet Listesi}
\begin{itemize}
    \item Kullanıcılar, katıldıkları tüm sohbetleri görebilmelidir.
    \item Sohbet listesi, son mesaj zamanına göre sıralanmalıdır.
\end{itemize}

\subsection{Mesajlaşma Gereksinimleri}
\textbf{FR-3.1: Mesaj Gönderme}
\begin{itemize}
    \item Kullanıcılar, sohbete metin mesajı gönderebilmelidir.
    \item Mesajlar gerçek zamanlı olarak alıcılara iletilebilmelidir.
    \item Mesaj gönderildi/iletildi/okundu durumları takip edilmelidir.
\end{itemize}

\textbf{FR-3.2: Mesaj Geçmişi}
\begin{itemize}
    \item Kullanıcılar, sohbetin mesaj geçmişini görüntüleyebilmelidir.
    \item Mesaj geçmişi sayfalama (pagination) ile yüklenebilmelidir.
\end{itemize}

\textbf{FR-3.3: Mesaj Silme}
\begin{itemize}
    \item Kullanıcılar, gönderdikleri mesajları silebilmelidir.
    \item Mesaj veritabanından fiziksel olarak silinmemelidir, soft delete uygulanmalıdır.
\end{itemize}

\subsection{Dosya Paylaşımı Gereksinimleri}
\textbf{FR-4.1: Dosya Yükleme}
\begin{itemize}
    \item Kullanıcılar, sohbete dosya yükleyebilmelidir.
    \item Maksimum dosya boyutu 50 MB olmalıdır.
    \item Dosya yüklendikten sonra bir URL oluşturulmalıdır.
\end{itemize}

\textbf{FR-4.2: Dosya İndirme}
\begin{itemize}
    \item Kullanıcılar, paylaşılan dosyaları indirebilmelidir.
    \item Dosya URL'leri güvenli olmalıdır.
\end{itemize}

\subsection{Bildirim Gereksinimleri}
\textbf{FR-5.1: Gerçek Zamanlı Bildirimler}
\begin{itemize}
    \item Yeni mesaj alındığında çevrimiçi kullanıcılara anında bildirim gönderilmelidir.
    \item Bildirimler WebSocket üzerinden iletilmelidir.
\end{itemize}

\textbf{FR-5.2: Çevrimdışı Bildirimler}
\begin{itemize}
    \item Çevrimdışı kullanıcılar için bildirimler kuyrukta saklanmalıdır.
\end{itemize}

\section{İşlevsel Olmayan Gereksinimler (Non-Functional Requirements)}

\subsection{Performans Gereksinimleri}
\textbf{NFR-1.1: Yanıt Süresi}
\begin{itemize}
    \item API isteklerinin \%95'i 500 ms içinde yanıtlanmalıdır.
    \item Mesaj gönderme işlemi 200 ms içinde tamamlanmalıdır.
    \item Gerçek zamanlı mesaj iletimi 100 ms içinde yapılmalıdır.
\end{itemize}

\textbf{NFR-1.2: İşlem Kapasitesi}
\begin{itemize}
    \item Sistem, saniyede en az 1000 mesaj işleyebilmelidir.
    \item Aynı anda 10,000 aktif kullanıcıyı destekleyebilmelidir.
\end{itemize}

\subsection{Ölçeklenebilirlik Gereksinimleri}
\textbf{NFR-2.1: Yatay Ölçeklenebilirlik}
\begin{itemize}
    \item Servisler bağımsız olarak ölçeklendirilebilmelidir.
    \item Her servis en az 3 örnekle çalışabilmelidir.
\end{itemize}

\textbf{NFR-2.2: Veri Ölçeklenebilirliği}
\begin{itemize}
    \item Veritabanı şeması milyonlarca mesajı destekleyebilmelidir.
\end{itemize}

\subsection{Güvenlik Gereksinimleri}
\textbf{NFR-3.1: Kimlik Doğrulama ve Yetkilendirme}
\begin{itemize}
    \item Tüm API istekleri JWT token ile doğrulanmalıdır.
    \item Şifreler hash'lenmiş (bcrypt) olarak saklanmalıdır.
\end{itemize}

\textbf{NFR-3.2: Veri Güvenliği}
\begin{itemize}
    \item API iletişimi HTTPS üzerinden yapılmalıdır.
    \item Dosya erişim URL'leri zaman sınırlı olmalıdır.
\end{itemize}

\textbf{NFR-3.3: Güvenlik Kontrolleri}
\begin{itemize}
    \item Rate limiting uygulanmalıdır (kullanıcı başına dakikada maksimum 60 istek).
\end{itemize}

\subsection{Dayanıklılık (Reliability) Gereksinimleri}
\textbf{NFR-5.1: Sistem Uptime}
\begin{itemize}
    \item Sistem \%99.5 uptime sağlamalıdır.
\end{itemize}

\textbf{NFR-5.2: Hata Toleransı}
\begin{itemize}
    \item Tek bir servisin çökmesi tüm sistemi etkilememelidir.
    \item Circuit breaker pattern kullanılarak servis hataları yönetilmelidir.
\end{itemize}

\subsection{Bakım Gereksinimleri}
\textbf{NFR-6.1: Loglama}
\begin{itemize}
    \item Tüm servisler merkezi loglama sistemine log göndermelidir.
    \item ELK Stack kullanılmalıdır.
\end{itemize}

\textbf{NFR-6.2: İzleme (Monitoring)}
\begin{itemize}
    \item Sistem sağlığı metrikleri izlenmelidir.
    \item API yanıt süreleri ve hata oranları izlenmelidir.
\end{itemize}
